\SourcePage{761}{764}
\section{Inelastic collisions of heavy particles with atoms}

Expressed in terms of the particle speed, the condition for applicability of
the Born approximation to collisions of heavy particles with atoms remains
the same as for electrons: $v\gg v_0$.

This follows directly from the general condition (126.2) for applicability of
perturbation theory, $Ua_0/\hbar v\ll1$, upon observing that the particle mass
does not enter it at all, while $Ua_0/\hbar$ is of the order of the atomic
electron speed.

In the coordinate system in which the center of mass of the atom and particle
is at rest, the cross section is determined by the general formula (148.3)
(where $m$ must now mean the reduced mass of the particle and atom). It is
more convenient, however, to consider the collision in the coordinate system
in which the scattering atom is at rest before the collision. For this
purpose we start from (148.1). In the coordinate system in which the atom was
at rest before the collision, the argument of the delta-function expressing
energy conservation is
\begin{equation}
 \frac{p'^2}{2M}-\frac{p^2}{2M}+
 \frac{(\mathbf p'-\mathbf p)^2}{2M_a}+E_n-E_0,
 \tag{150.1}
\end{equation}
where $M$ is the mass of the incident particle and $M_a$ the mass of the atom;
the third term is the atom's recoil kinetic energy (which could be neglected
completely in an electron collision).

When a fast heavy particle collides with an atom, the change in the particle's
momentum is almost always small compared with its initial momentum. If this
condition holds, the atom's recoil energy may be neglected in the argument of
the delta-function. We then return exactly to formula (148.3), except that
$m$ must be replaced by the mass $M$ of the incident particle (not by the
reduced mass of the particle and atom). Since the momentum transfer is
assumed small compared with the initial momentum, we put $p\approx p'$.
Thus, for the cross section in the coordinate system in which the atom was at
rest before the collision,

\SourcePage{762}{765}
we obtain
\begin{equation}
 d\sigma_n=\frac{M^2}{4\pi^2\hbar^4}
 \left|\iint Ue^{-i\mathbf q\mathbf r}\psi_n^*\psi_0\,d\tau\,dV\right|^2d\Omega.
 \tag{150.2}
\end{equation}

To allow for a particle charge different from the electron charge, we write
$ze^2$ in place of $e^2$, where $ze$ is the charge of the incident particle.
The general formula for inelastic scattering, written in the form (148.9),
\begin{equation}
 d\sigma_n=8\pi\left(\frac{ze^2}{\hbar v}\right)^2
 \left|\left\langle n\left|\sum_a e^{-i\mathbf q\mathbf r_a}\right|0\right\rangle\right|^2
 \frac{dq}{q^3},
 \tag{150.3}
\end{equation}
contains no particle mass. It follows that all formulas derived from it remain
applicable to collisions of heavy particles as well, provided those formulas
are expressed in terms of $v$ and $q$.

It is readily seen how formulas expressed in terms of the scattering angle
$\vartheta$ (the deflection angle of the heavy particle colliding with the
atom) must be modified. First observe that in an inelastic collision of a
heavy particle, $\vartheta$ is always small. Indeed, when the momentum
transfer is large compared with atomic-electron momenta, an inelastic
collision with an atom may be regarded as an elastic collision with free
electrons; but when a heavy particle collides with a light one (an electron),
the heavy particle is hardly deflected. In other words, the momentum transfer
from the heavy particle to the atom is small compared with the particle's
initial momentum (the exception is elastic scattering through large angles,
which is, however, extremely improbable).

Consequently, throughout the angular range one may put
\begin{equation}
 q=\frac1\hbar\sqrt{\left(\frac{E_n-E_0}{v}\right)^2+(Mv\vartheta)^2},
 \tag{150.4}
\end{equation}
which in practice reduces to
\begin{equation}
 q\hbar\approx Mv\vartheta
 \tag{150.5}
\end{equation}
except at the very smallest angles. On the other hand, when considering
electron collisions with an atom, we wrote (for small angles)
\[
 q=\frac1\hbar\sqrt{\left(\frac{E_n-E_0}{v}\right)^2+(mv\vartheta)^2}.
\]
Comparison of the two expressions shows that the formulas obtained for
electron collisions with atoms,

\SourcePage{763}{766}
when expressed in terms of the speed and deflection angle, are converted into
formulas for heavy-particle collisions by making the replacement everywhere
(including in the solid-angle element
$d\Omega=2\pi\sin\vartheta\,d\vartheta\approx
2\pi\vartheta\,d\vartheta$)
\begin{equation}
 \vartheta\longrightarrow\frac{M}{m}\vartheta,
 \tag{150.6}
\end{equation}
at the same incident-particle speed $v$. Qualitatively, this means that the
entire small-angle scattering pattern is narrowed by the factor $m/M$ (at a
specified speed).

The rules obtained also apply to elastic scattering at small angles. Applying
the transformation (150.6) to (139.4) with $\vartheta\ll1$, we obtain the
cross section
\begin{equation}
 d\sigma_e=8\pi\left(\frac{ze^2}{Mv^2}\right)^2
 \left[Z-F\left(\frac{Mv\vartheta}{\hbar}\right)\right]^2
 \frac{d\vartheta}{\vartheta^3}.
 \tag{150.7}
\end{equation}
Elastic scattering of heavy particles through angles $\vartheta\sim1$, on the
other hand, reduces to Rutherford scattering by the atomic nucleus.

Inelastic scattering with ionization of the atom at large momentum transfer
requires separate consideration. Unlike electron-impact ionization, there
are of course no exchange effects here. A characteristic feature of heavy
particles is that a large momentum transfer ($qa_0\gg1$) does not at all imply
deflection through a large angle; $\vartheta$ always remains small. The
ionization cross section for emission of an electron with energy between
$\varepsilon$ and $\varepsilon+d\varepsilon$ follows directly from (148.25),
which we write as
\[
 d\sigma_r=8\pi\left(\frac{ze^2}{\hbar v}\right)^2Z\frac{dq}{q^3},
\]
and in which we put $\hbar^2q^2/2m=\varepsilon$ (the entire momentum
$\hbar\mathbf q$ is transferred to one atomic electron). This gives
\begin{equation}
 d\sigma_r=\frac{2\pi Zz^2e^4}{mv^2}
 \frac{d\varepsilon}{\varepsilon^2}.
 \tag{150.8}
\end{equation}

For collisions of heavy particles with atoms, integral effective cross
sections and stopping are of particular interest.

The total inelastic-scattering cross section is determined by the previous
formula (148.26). The total effective stopping is obtained by substituting in
(149.12), for $q_1$, the greatest possible momentum transfer $q_{\max}$. The
latter is readily expressed in terms of the particle speed as follows. Since
$\hbar q_{\max}$ is still small compared with the particle's initial momentum
$Mv$,

\SourcePage{764}{767}
the change in its energy is related to the change in momentum by
$\Delta E=\mathbf v\cdot\hbar\mathbf q$.

On the other hand, at large momentum transfer essentially all this energy is
transferred to one atomic electron, so that we may write
\[
 \varepsilon=\frac{\hbar^2q^2}{2m}
 =\hbar\mathbf v\mathbf q\leqslant\hbar vq.
\]
It follows that $\hbar q\leqslant2mv$, that is,
\begin{equation}
 \hbar q_{\max}=2mv,\qquad \varepsilon_{\max}=2mv^2.
 \tag{150.9}
\end{equation}
Note that the greatest particle deflection angle in inelastic scattering is
\[
 \vartheta_{\max}=\frac{\hbar q_{\max}}{Mv}=\frac{2m}{M}.
\]
Substitution of (150.9) in (149.12) gives the total effective stopping of a
heavy particle:
\begin{equation}
 \kappa=\frac{4\pi Zz^2e^4}{mv^2}\ln\frac{2mv^2}{I}.
 \tag{150.10}
\end{equation}
